\documentclass[10pt,journal,compsoc,twocolumn]{IEEEtran}
\usepackage[utf8]{inputenc}
\usepackage{amsmath,amssymb,amsfonts}
\usepackage{graphicx}
\usepackage{booktabs}
\usepackage{cite}
\usepackage{hyperref}
\usepackage{microtype}

\begin{document}

\title{Task B2: Multi-Modal Wine Critic-Score Prediction via Tabular-Text Feature Fusion and Ablation}

\author{
    \IEEEauthorblockN{Muhammad Abdullah}\\
    \IEEEauthorblockA{Department of Intelligent Systems, The University of Lahore, Pakistan\\
    Roll No: 70178208 \quad|\quad Section: BSRAI-5C\\
    GitHub: \url{https://github.com/MuhammadaAbdullah70178208/AI101-ML-Assignment-70178208}\\
    Instructor: Dr. Muhammad Awais, AI Architect}
}

\maketitle

\begin{abstract}
Evaluating wine quality typically involves combining objective market attributes (retail price, origin, cultivar) with nuanced sensory notes recorded by certified sommeliers. This investigation addresses Task B2 of the AI-101 curriculum, evaluating multi-modal regression methodologies across tabular and unstructured textual modalities on a 5,000-sample Wine Magazine benchmark. We implement a rigorous modality ablation experiment isolating tabular metadata, n-gram TF-IDF textual features, and late concatenated representations. While a Tabular-only Ridge regressor achieves an RMSE of 1.464 points ($R^2 = 0.8154$) and a Text-only regressor achieves 1.557 points ($R^2 = 0.7912$), combining both representations yields a test RMSE of 1.121 points ($R^2 = 0.8917$). This establishes an empirical error reduction of 23.4\% over tabular attributes alone, demonstrating that unstructured sensory descriptions provide critical orthogonal signals for continuous appraisal.
\end{abstract}

\section{Introduction}
Quantifying organoleptic wine quality is economically pivotal for viticulture valuation, competitive pricing, and consumer cellar curation. Wine reviews published by major rating institutions follow a dual structure: structured categorical and numeric metadata (grape variety, country of origin, bottle price) paired with subjective free-form tasting notes. Modeling this problem requires multi-modal machine learning capable of uniting dense tabular matrices with sparse, high-dimensional linguistic representations while preventing data leakage across disjoint splits.

\section{Problem Statement}
We frame critic score estimation as a supervised multi-modal continuous regression problem:
\begin{equation}
f: \mathcal{X}_{\text{tab}} \times \mathcal{X}_{\text{text}} \to \mathbb{R}
\end{equation}
where $\mathbf{x}_{\text{tab}} = [\log(1 + p), \mathbf{c}_{\text{ohe}}, \mathbf{v}_{\text{ohe}}] \in \mathbb{R}^{21}$ encompasses transformed bottle price $p$ and one-hot encoded country and grape variety vectors, and $\mathbf{x}_{\text{text}} \in \mathbb{R}^{3000}$ denotes a TF-IDF n-gram vector extracted from sommelier tasting reviews. The continuous target $y \in [80, 100]$ represents the 100-point critic score. Primary evaluation metrics are Mean Absolute Error (MAE), Root Mean Squared Error (RMSE), and the coefficient of determination ($R^2$).

\section{Dataset and Preprocessing}
We utilize the Wine Magazine benchmark dataset (5,000 deduplicated records). The dataset features a right-skewed price distribution with 7.0\% missing observations.

\begin{table}[htbp]
\centering
\caption{WineMag Multi-Modal Dataset Characteristics}
\label{tab:winemag_stats}
\begin{tabular}{ll}
\toprule
\textbf{Attribute} & \textbf{Specification} \\
\midrule
Total Samples ($N$) & 5,000 reviews \\
Target Variable & Points ($80$ to $100$, $\mu = 88.44 \pm 3.41$) \\
Tabular Attributes & Country (10 classes), Variety (10 classes), Price \\
Text Modality & Tasting notes ($15$ to $40$ words per review) \\
Missing Price Rate & 7.0\% (imputed with training-set median) \\
Splitting Protocol & Disjoint 70\% Train ($3,500$) / 15\% Val ($750$) \\
                   & / 15\% Test ($750$) \\
\bottomrule
\end{tabular}
\end{table}

\begin{figure}[htbp]
\centering
\includegraphics[width=\columnwidth]{figures/b2_eda_price_points_distribution.png}
\caption{Task B2 EDA: (Left) 100-point critic score distribution; (Right) Positive log-linear correlation between bottle price and critic rating ($r = 0.61$).}
\label{fig:eda}
\end{figure}

To prevent data leakage in compliance with Module~02, all preprocessing transformations are computed strictly over the training partition:
\begin{enumerate}
    \item \textbf{Price Transformation:} $p' = \log(1 + \text{price})$, imputing missing entries with the training median ($\text{med} = \$28.50$).
    \item \textbf{Categorical Encoding:} Sparse one-hot encoding for country and variety with unseen-category handling.
    \item \textbf{Text Vectorization:} Sub-linear TF-IDF vectorization with $(1, 2)$-gram collocations, English stop-word filtering, and top 3,000 vocabulary selection.
\end{enumerate}

\section{Experimental Architectures and Ablation}
We benchmark five modeling regimes to evaluate modality contributions:
\begin{enumerate}
    \item \textbf{Baseline Mean Regressor:} Dummy regressor predicting $\bar{y}_{\text{train}}$.
    \item \textbf{Tabular-Only Ridge:} $\ell_2$-regularized regression trained strictly on standardized tabular vectors ($\alpha = 10.0$).
    \item \textbf{Text-Only Ridge:} $\ell_2$-regularized linear model trained exclusively on the 3,000-dimensional TF-IDF matrix ($\alpha = 1.0$).
    \item \textbf{Combined Ridge (Multi-Modal):} Ridge regression trained on the late concatenated feature space $\mathbf{x}_{\text{comb}} = [\mathbf{x}_{\text{tab}}, \mathbf{x}_{\text{text}}] \in \mathbb{R}^{3021}$ with $\alpha = 5.0$.
    \item \textbf{Combined Gradient Boosting:} Non-linear decision tree ensemble (100 estimators, depth 4) trained on concatenated features.
\end{enumerate}

\section{Results and Performance Analysis}
Test performance on 750 unseen evaluation instances is presented in Table~\ref{tab:results}.

\begin{table}[htbp]
\centering
\caption{Quantitative Performance on Unseen Test Partition ($N=750$)}
\label{tab:results}
\begin{tabular}{lccc}
\toprule
\textbf{Model Architecture} & \textbf{MAE} & \textbf{RMSE} & \textbf{$R^2$} \\
\midrule
Baseline Mean Regressor      & 2.731 & 3.407 & -0.0002 \\
Tabular-Only Ridge (Metadata)& 1.107 & 1.464 & 0.8154 \\
Text-Only Ridge (TF-IDF)     & 1.319 & 1.557 & 0.7912 \\
\textbf{Combined Ridge}      & \textbf{0.897} & \textbf{1.121} & \textbf{0.8917} \\
Combined Gradient Boosting   & 0.895 & 1.186 & 0.8789 \\
\bottomrule
\end{tabular}
\end{table}

\begin{figure}[htbp]
\centering
\includegraphics[width=\columnwidth]{figures/b2_ablation_and_residuals.png}
\caption{Task B2 Diagnostics: (Left) Modality ablation bar chart illustrating RMSE reduction across architectures; (Right) Homoscedastic residual error scatter for Combined Ridge.}
\label{fig:ablation_res}
\end{figure}

\section{Modality Ablation and Discussion}
The experimental results demonstrate the following:
\begin{enumerate}
    \item \textbf{Synergistic Modality Fusion:} The multi-modal Combined Ridge regressor yields a test RMSE of 1.121 points and an $R^2$ of 0.8917. This represents a 23.4\% reduction in RMSE compared to tabular features alone and a 28.0\% reduction over text alone.
    \item \textbf{Signal Complementarity:} While bottle price provides a strong macroscopic indicator of cellar tier ($r=0.61$), sommelier tasting text captures nuanced defects (``astringent'', ``medicinal'', ``flabby'') and premium indicators (``graphite'', ``opulent brioche'', ``majestic persistence'') necessary to fine-tune final ratings.
    \item \textbf{Residual Diagnostics:} The residual scatter plot in Fig.~\ref{fig:ablation_res} exhibits uniform, zero-centered dispersion across the entire dynamic range ($80$ to $98$ points), validating the stability of $\ell_2$ shrinkage against over-fitting to rare adjectives.
\end{enumerate}

\section{Conclusion}
This study established that late feature fusion combining structured wine metadata with TF-IDF text features achieves superior predictive precision (RMSE 1.121, $R^2 = 0.8917$). The modality ablation confirms that linguistic reviews provide substantial complementary variance over economic metadata alone.

\begin{thebibliography}{1}
\bibitem{thoutt2017wine}
Z.~Thoutt, ``Wine Magazine Dataset: Points and price ratings with sommelier descriptions,'' \emph{Kaggle Repository}, 2017.
\bibitem{mitchell1997machine}
T.~M.~Mitchell, \emph{Machine Learning}, McGraw-Hill, New York, 1997.
\bibitem{awais2026introai}
M.~Awais, \emph{AI-101: Introduction to Artificial Intelligence (Course Slides: Modules 01 \& 02)}, Dept. of Intelligent Systems, The University of Lahore, 2026.
\end{thebibliography}

\end{document}
